\IfFileExists{stacks-project.cls}{%
\documentclass{stacks-project}
}{%
\documentclass{amsart}
}

% For dealing with references we use the comment environment
\usepackage{verbatim}
\newenvironment{reference}{\comment}{\endcomment}
%\newenvironment{reference}{}{}
\newenvironment{slogan}{\comment}{\endcomment}
\newenvironment{history}{\comment}{\endcomment}

% For commutative diagrams we use Xy-pic
\usepackage[all]{xy}

% We use 2cell for 2-commutative diagrams.
\xyoption{2cell}
\UseAllTwocells

% We use multicol for the list of chapters between chapters
\usepackage{multicol}

% This is generally recommended for better output
\usepackage{lmodern}
\usepackage[T1]{fontenc}

% For cross-file-references
\usepackage{xr-hyper}

% Package for hypertext links:
\usepackage{hyperref}

% For any local file, say "hello.tex" you want to link to please
% use \externaldocument[hello-]{hello}
\externaldocument[introduction-]{introduction}
\externaldocument[conventions-]{conventions}
\externaldocument[sets-]{sets}
\externaldocument[categories-]{categories}
\externaldocument[topology-]{topology}
\externaldocument[sheaves-]{sheaves}
\externaldocument[sites-]{sites}
\externaldocument[stacks-]{stacks}
\externaldocument[fields-]{fields}
\externaldocument[algebra-]{algebra}
\externaldocument[brauer-]{brauer}
\externaldocument[homology-]{homology}
\externaldocument[derived-]{derived}
\externaldocument[simplicial-]{simplicial}
\externaldocument[more-algebra-]{more-algebra}
\externaldocument[smoothing-]{smoothing}
\externaldocument[modules-]{modules}
\externaldocument[sites-modules-]{sites-modules}
\externaldocument[injectives-]{injectives}
\externaldocument[cohomology-]{cohomology}
\externaldocument[sites-cohomology-]{sites-cohomology}
\externaldocument[dga-]{dga}
\externaldocument[dpa-]{dpa}
\externaldocument[sdga-]{sdga}
\externaldocument[hypercovering-]{hypercovering}
\externaldocument[schemes-]{schemes}
\externaldocument[constructions-]{constructions}
\externaldocument[properties-]{properties}
\externaldocument[morphisms-]{morphisms}
\externaldocument[coherent-]{coherent}
\externaldocument[divisors-]{divisors}
\externaldocument[limits-]{limits}
\externaldocument[varieties-]{varieties}
\externaldocument[topologies-]{topologies}
\externaldocument[descent-]{descent}
\externaldocument[perfect-]{perfect}
\externaldocument[more-morphisms-]{more-morphisms}
\externaldocument[flat-]{flat}
\externaldocument[groupoids-]{groupoids}
\externaldocument[more-groupoids-]{more-groupoids}
\externaldocument[etale-]{etale}
\externaldocument[chow-]{chow}
\externaldocument[intersection-]{intersection}
\externaldocument[pic-]{pic}
\externaldocument[weil-]{weil}
\externaldocument[adequate-]{adequate}
\externaldocument[dualizing-]{dualizing}
\externaldocument[duality-]{duality}
\externaldocument[discriminant-]{discriminant}
\externaldocument[derham-]{derham}
\externaldocument[local-cohomology-]{local-cohomology}
\externaldocument[algebraization-]{algebraization}
\externaldocument[curves-]{curves}
\externaldocument[resolve-]{resolve}
\externaldocument[models-]{models}
\externaldocument[functors-]{functors}
\externaldocument[equiv-]{equiv}
\externaldocument[pione-]{pione}
\externaldocument[etale-cohomology-]{etale-cohomology}
\externaldocument[proetale-]{proetale}
\externaldocument[relative-cycles-]{relative-cycles}
\externaldocument[more-etale-]{more-etale}
\externaldocument[trace-]{trace}
\externaldocument[crystalline-]{crystalline}
\externaldocument[spaces-]{spaces}
\externaldocument[spaces-properties-]{spaces-properties}
\externaldocument[spaces-morphisms-]{spaces-morphisms}
\externaldocument[decent-spaces-]{decent-spaces}
\externaldocument[spaces-cohomology-]{spaces-cohomology}
\externaldocument[spaces-limits-]{spaces-limits}
\externaldocument[spaces-divisors-]{spaces-divisors}
\externaldocument[spaces-over-fields-]{spaces-over-fields}
\externaldocument[spaces-topologies-]{spaces-topologies}
\externaldocument[spaces-descent-]{spaces-descent}
\externaldocument[spaces-perfect-]{spaces-perfect}
\externaldocument[spaces-more-morphisms-]{spaces-more-morphisms}
\externaldocument[spaces-flat-]{spaces-flat}
\externaldocument[spaces-groupoids-]{spaces-groupoids}
\externaldocument[spaces-more-groupoids-]{spaces-more-groupoids}
\externaldocument[bootstrap-]{bootstrap}
\externaldocument[spaces-pushouts-]{spaces-pushouts}
\externaldocument[spaces-chow-]{spaces-chow}
\externaldocument[groupoids-quotients-]{groupoids-quotients}
\externaldocument[spaces-more-cohomology-]{spaces-more-cohomology}
\externaldocument[spaces-simplicial-]{spaces-simplicial}
\externaldocument[spaces-duality-]{spaces-duality}
\externaldocument[formal-spaces-]{formal-spaces}
\externaldocument[restricted-]{restricted}
\externaldocument[spaces-resolve-]{spaces-resolve}
\externaldocument[formal-defos-]{formal-defos}
\externaldocument[defos-]{defos}
\externaldocument[cotangent-]{cotangent}
\externaldocument[examples-defos-]{examples-defos}
\externaldocument[algebraic-]{algebraic}
\externaldocument[examples-stacks-]{examples-stacks}
\externaldocument[stacks-sheaves-]{stacks-sheaves}
\externaldocument[criteria-]{criteria}
\externaldocument[artin-]{artin}
\externaldocument[quot-]{quot}
\externaldocument[stacks-properties-]{stacks-properties}
\externaldocument[stacks-morphisms-]{stacks-morphisms}
\externaldocument[stacks-limits-]{stacks-limits}
\externaldocument[stacks-cohomology-]{stacks-cohomology}
\externaldocument[stacks-perfect-]{stacks-perfect}
\externaldocument[stacks-introduction-]{stacks-introduction}
\externaldocument[stacks-more-morphisms-]{stacks-more-morphisms}
\externaldocument[stacks-geometry-]{stacks-geometry}
\externaldocument[moduli-]{moduli}
\externaldocument[moduli-curves-]{moduli-curves}
\externaldocument[examples-]{examples}
\externaldocument[exercises-]{exercises}
\externaldocument[guide-]{guide}
\externaldocument[desirables-]{desirables}
\externaldocument[coding-]{coding}
\externaldocument[obsolete-]{obsolete}
\externaldocument[fdl-]{fdl}
\externaldocument[index-]{index}

% Theorem environments.
%
\theoremstyle{plain}
\newtheorem{theorem}[subsection]{Theorem}
\newtheorem{proposition}[subsection]{Proposition}
\newtheorem{lemma}[subsection]{Lemma}

\theoremstyle{definition}
\newtheorem{definition}[subsection]{Definition}
\newtheorem{example}[subsection]{Example}
\newtheorem{exercise}[subsection]{Exercise}
\newtheorem{situation}[subsection]{Situation}

\theoremstyle{remark}
\newtheorem{remark}[subsection]{Remark}
\newtheorem{remarks}[subsection]{Remarks}

\numberwithin{equation}{subsection}

% Macros
%
\def\lim{\mathop{\mathrm{lim}}\nolimits}
\def\colim{\mathop{\mathrm{colim}}\nolimits}
\def\Spec{\mathop{\mathrm{Spec}}}
\def\Hom{\mathop{\mathrm{Hom}}\nolimits}
\def\Ext{\mathop{\mathrm{Ext}}\nolimits}
\def\SheafHom{\mathop{\mathcal{H}\!\mathit{om}}\nolimits}
\def\SheafExt{\mathop{\mathcal{E}\!\mathit{xt}}\nolimits}
\def\Sch{\mathit{Sch}}
\def\Mor{\mathop{\mathrm{Mor}}\nolimits}
\def\Ob{\mathop{\mathrm{Ob}}\nolimits}
\def\Sh{\mathop{\mathit{Sh}}\nolimits}
\def\NL{\mathop{N\!L}\nolimits}
\def\CH{\mathop{\mathrm{CH}}\nolimits}
\def\proetale{{pro\text{-}\acute{e}tale}}
\def\etale{{\acute{e}tale}}
\def\QCoh{\mathit{QCoh}}
\def\Ker{\mathop{\mathrm{Ker}}}
\def\Im{\mathop{\mathrm{Im}}}
\def\Coker{\mathop{\mathrm{Coker}}}
\def\Coim{\mathop{\mathrm{Coim}}}

% Boxtimes
%
\DeclareMathSymbol{\boxtimes}{\mathbin}{AMSa}{"02}

%
% Macros for moduli stacks/spaces
%
\def\QCohstack{\mathcal{QC}\!\mathit{oh}}
\def\Cohstack{\mathcal{C}\!\mathit{oh}}
\def\Spacesstack{\mathcal{S}\!\mathit{paces}}
\def\Quotfunctor{\mathrm{Quot}}
\def\Hilbfunctor{\mathrm{Hilb}}
\def\Curvesstack{\mathcal{C}\!\mathit{urves}}
\def\Polarizedstack{\mathcal{P}\!\mathit{olarized}}
\def\Complexesstack{\mathcal{C}\!\mathit{omplexes}}
% \Pic is the operator that assigns to X its picard group, usage \Pic(X)
% \Picardstack_{X/B} denotes the Picard stack of X over B
% \Picardfunctor_{X/B} denotes the Picard functor of X over B
\def\Pic{\mathop{\mathrm{Pic}}\nolimits}
\def\Picardstack{\mathcal{P}\!\mathit{ic}}
\def\Picardfunctor{\mathrm{Pic}}
\def\Deformationcategory{\mathcal{D}\!\mathit{ef}}
\usepackage{cleveref}
\crefname{originalproposition}{original proposition}{original propositions}
\Crefname{originalproposition}{Original proposition}{Original propositions}
\makeatletter
\newlabel{proposition-pic-functor}{{0D2C}{source}{}{sourcepointer-1}{}}
\newlabel{proposition-pic-functor@cref}{{[originalproposition][0][]0D2C}{[1][0][]source}{}{}{}}
\makeatother
\crefname{originallemma}{original lemma}{original lemmas}
\Crefname{originallemma}{Original lemma}{Original lemmas}
\makeatletter
\newlabel{more-morphisms-lemma-separated-locally-quasi-finite-morphisms-fppf-descend}{{02W8}{source}{}{sourcepointer-2}{}}
\newlabel{more-morphisms-lemma-separated-locally-quasi-finite-morphisms-fppf-descend@cref}{{[originallemma][0][]02W8}{[1][0][]source}{}{}{}}
\makeatother
\crefname{originallemma}{original lemma}{original lemmas}
\Crefname{originallemma}{Original lemma}{Original lemmas}
\makeatletter
\newlabel{descent-lemma-descending-fppf-property-immersion}{{02YM}{source}{}{sourcepointer-3}{}}
\newlabel{descent-lemma-descending-fppf-property-immersion@cref}{{[originallemma][0][]02YM}{[1][0][]source}{}{}{}}
\makeatother
\crefname{originallemma}{original lemma}{original lemmas}
\Crefname{originallemma}{Original lemma}{Original lemmas}
\makeatletter
\newlabel{lemma-flat-geometrically-connected-fibres}{{0D25}{source}{}{sourcepointer-4}{}}
\newlabel{lemma-flat-geometrically-connected-fibres@cref}{{[originallemma][0][]0D25}{[1][0][]source}{}{}{}}
\makeatother
\crefname{originallemma}{original lemma}{original lemmas}
\Crefname{originallemma}{Original lemma}{Original lemmas}
\makeatletter
\newlabel{spaces-more-morphisms-lemma-diagonal-picard-flat-proper}{{0D23}{source}{}{sourcepointer-5}{}}
\newlabel{spaces-more-morphisms-lemma-diagonal-picard-flat-proper@cref}{{[originallemma][0][]0D23}{[1][0][]source}{}{}{}}
\makeatother
\title{The diagonal of the Picard functor is representable by immersions}
\author{The Stacks Project Authors}
\date{stacks-project:quot.tex:local-snapshot-20261012; GNU FDL 1.2 or later}
\begin{document}
\maketitle
\setcounter{subsection}{1}
\noindent\textit{Editorial note.} All delivered mathematical text is the author's own, copied from the pinned Stacks Project chapter file quot.tex; no statement or proof step was rewritten, reordered or removed. The counted span is the author environment and its own proof environment. References inside the span are handed to the assembler as navigation stubs that carry the canonical tag of the cited tag only; whether each of them is load-bearing is recorded on the stub itself.\par\medskip
\noindent\textit{Allowed standard prerequisites (editorial).}\begin{itemize}
\item Proposition proposition-pic-functor (canonical tag 0D2C). cited at statement line 3275. Assumptions and notation: the Picard functor is an fppf sheaf.
\item Lemma more-morphisms-lemma-separated-locally-quasi-finite-morphisms-fppf-descend (canonical tag 02W8). cited at proof lines 3299-3301. Descent of the locally quasi-finite Z\_i -> T\_i to Z -> T.
\item Lemma descent-lemma-descending-fppf-property-immersion (canonical tag 02YM). cited at proof line 3304. Z -> T is an immersion.
\item Lemma lemma-flat-geometrically-connected-fibres (canonical tag 0D25). cited at proof line 3316. Characterizes Z: factoring through Z iff L\_\{X\_\{T'\}\} is a pullback from T'.
\item Lemma spaces-more-morphisms-lemma-diagonal-picard-flat-proper (canonical tag 0D23). cited at proof lines 3317-3319. Existence of the locally closed Z in the invertible-module case.
\end{itemize}\medskip
\ifdefined\hypertarget\hypertarget{sourcepointer-1}{}\fi
\noindent\textit{Source reference.} Original proposition 0D2C: sources/stacks-chow/tags.txt, source lines 16377--16377. Cited at statement line 3275: assumptions and notation inherited from this proposition; load-bearing.\par\medskip
\ifdefined\hypertarget\hypertarget{sourcepointer-2}{}\fi
\noindent\textit{Source reference.} Original lemma 02W8: sources/stacks-chow/tags.txt, source lines 3795--3795. Cited at proof lines 3299-3301: Z -> T obtained by descending the locally quasi-finite Z\_i -> T\_i; load-bearing.\par\medskip
\ifdefined\hypertarget\hypertarget{sourcepointer-3}{}\fi
\noindent\textit{Source reference.} Original lemma 02YM: sources/stacks-chow/tags.txt, source lines 3881--3881. Cited at proof line 3304: Z -> T is an immersion; load-bearing.\par\medskip
\ifdefined\hypertarget\hypertarget{sourcepointer-4}{}\fi
\noindent\textit{Source reference.} Original lemma 0D25: sources/stacks-chow/tags.txt, source lines 16370--16370. Cited at proof line 3316: T' -> T factors through Z iff L\_\{X\_\{T'\}\} is the pullback of an invertible sheaf on T'; load-bearing.\par\medskip
\ifdefined\hypertarget\hypertarget{sourcepointer-5}{}\fi
\noindent\textit{Source reference.} Original lemma 0D23: sources/stacks-chow/tags.txt, source lines 16368--16368. Cited at proof lines 3317-3319: existence of Z in the invertible-module case; load-bearing.\par\medskip
\begin{lemma}
\label{lemma-diagonal-pic}
With assumptions and notation as in Proposition \ref{proposition-pic-functor}.
Then the diagonal
$\Picardfunctor_{X/B} \to \Picardfunctor_{X/B} \times_B \Picardfunctor_{X/B}$
is representable by immersions. In other words, $\Picardfunctor_{X/B} \to B$
is locally separated.
\end{lemma}

\begin{proof}
Let $T$ be a scheme over $B$ and let $s, t \in \Picardfunctor_{X/B}(T)$.
We want to show that there exists a locally closed subscheme $Z \subset T$
such that $s|_Z = t|_Z$ and such that a morphism $T' \to T$ factors
through $Z$ if and only if $s|_{T'} = t|_{T'}$.

\medskip\noindent
We first reduce the general problem to the case where $s$ and $t$ come
from invertible modules on $X_T$. We suggest the reader skip this step.
Choose an fppf covering $\{T_i \to T\}_{i \in I}$ such that
$s|_{T_i}$ and $t|_{T_i}$ come from $\Pic(X_{T_i})$ for all $i$.
Suppose that we can show the result for all the pairs
$s|_{T_i}, t|_{T_i}$. Then we obtain locally closed subschemes
$Z_i \subset T_i$ with the desired universal property.
It follows that $Z_i$ and $Z_j$ have the same scheme theoretic
inverse image in $T_i \times_T T_j$.
This determines a descend datum on $Z_i/T_i$.
Since $Z_i \to T_i$ is locally quasi-finite, it follows from
More on Morphisms, Lemma
\ref{more-morphisms-lemma-separated-locally-quasi-finite-morphisms-fppf-descend}
that we obtain a locally quasi-finite morphism $Z \to T$
recovering $Z_i \to T_i$ by base change. Then $Z \to T$ is an immersion
by Descent, Lemma \ref{descent-lemma-descending-fppf-property-immersion}.
Finally, because $\Picardfunctor_{X/B}$ is an fppf sheaf, we conclude
that $s|_Z = t|_Z$ and that $Z$ satisfies the universal property
mentioned above.

\medskip\noindent
Assume $s$ and $t$ come from invertible modules $\mathcal{V}$, $\mathcal{W}$
on $X_T$.
Set $\mathcal{L} = \mathcal{V} \otimes \mathcal{W}^{\otimes -1}$
We are looking for a locally closed subscheme $Z$ of $T$
such that $T' \to T$ factors through $Z$ if and only if $\mathcal{L}_{X_{T'}}$
is the pullback of an invertible sheaf on $T'$, see
Lemma \ref{lemma-flat-geometrically-connected-fibres}.
Hence the existence of $Z$ follows from
More on Morphisms of Spaces, Lemma
\ref{spaces-more-morphisms-lemma-diagonal-picard-flat-proper}.
\end{proof}

\end{document}
